% !TEX root = ../supporting_information.tex

\SISection{Physics-based upscaling of local fault properties}
\label{sec:si-local-fault-property-upscaling}

\SISubsection{Exact realization replay}

Every selected library row stores the original realization identifier,
PREDICT random seed, checkpoint hash, code commit, and method-configuration
hash. Replay uses this provenance record to reconstruct the selected
fine-scale parent units, clay-smear occupancy, porosity, and permeability
fields without archiving every original fine-scale realization. The replayed
effective permeability must agree with the frozen library value within the
validated numerical-equivalence tolerance, while the realization identifier,
seed, code and configuration hashes, and discrete material architecture must
match exactly. A failed check stops downstream upscaling for that fault
element rather than silently substituting another realization.

\todo{Report the final numerical-equivalence tolerance, checkpoint schema,
software versions, and replay-validation coverage across all selected
realizations.}

\SISubsection{Effective permeability and porosity}

Effective permeability is computed with the flow-based PREDICT upscaling
calculation, which applies pressure differences to the replayed fine-scale
fault core and obtains the directional flux responses. Effective porosity is
the pore-volume-weighted average of the fine-cell porosities. Permeability and
porosity are linked to the same replayed realization and are never assembled
from different source rows.

PREDICT reports the permeability tensor in its local fault coordinate system.
Before reservoir export, the tensor is mapped to reservoir-grid coordinates
using
\begin{equation}
  \mathbf{K}_{\mathrm{res}} =
  \mathbf{R}\mathbf{K}_{\mathrm{local}}\mathbf{R}^{\mathsf T},
  \label{eq:si-tensor-rotation}
\end{equation}
where $\mathbf{R}$ contains the local fault-normal, strike, and dip directions
expressed in reservoir coordinates. Exported files retain the local tensor,
rotation metadata, transformed reservoir tensor, component order, and units.
This explicit transformation prevents local $k_{xx}$, $k_{yy}$, and $k_{zz}$
from being interpreted directly as reservoir-grid components.

\todo{State the final local-axis convention, give the explicit rotation
matrix, document the permeability-unit conversion, and report tensor-
transformation validation against the downstream importer.}

\SISubsection{Capillary-pressure upscaling}

Capillary-pressure curves are upscaled by connected invasion percolation on
the replayed fine grid. Local entry pressures are obtained by scaling the
reference material relations with the corresponding fine-cell permeability
and porosity. At each imposed invasion threshold, the connected invaded pore
volume determines the bulk gas saturation. Continuing the invasion sequence
therefore produces both the effective drainage curve and its native maximum
gas-saturation endpoint; the effective irreducible water saturation is one
minus that endpoint.

The production reservoir input retains a separate full-slice capillary-
pressure curve for every one of the 522 fault elements. This representation
preserves along-strike variation in entry pressure, curve shape, and
saturation endpoint. An entry-pressure-branch-medoid representation is
retained only as a documented sensitivity option and is not used for the
production cases described in the main text.

\todo{Provide the reference material curves, Leverett-scaling relation,
invasion boundary condition, connectivity convention, threshold sequence,
and numerical interpolation used for reservoir export.}

\SISubsection{Relative-permeability upscaling}

Dynamic relative-permeability upscaling solves the validated fine-grid flow
problem for an actual replayed slice. For each throw window, the representative
slice is selected as the medoid of the effective irreducible-water-saturation
values obtained from the 87 capillary-pressure upscalings. The resulting gas
and water relative-permeability curve shapes are then mapped to the physical
saturation endpoints of the associated capillary-pressure regions. This
reduces the number of dynamic calculations while preserving the slice-
specific capillary endpoints required for reservoir simulation.

The robust dynamic solver uses internally controlled time stepping and the
validated AMGCL linear solver where supported. The workflow disallows
along-strike grid collapse and representative-slice substitution during the
capillary-pressure calculations. Relative-permeability curves must be
monotonic in the expected direction, remain within $[0,1]$, and terminate at
endpoints consistent with their assigned capillary-pressure regions.

\todo{Describe the fine-grid flow equations, phase properties, boundary and
initial conditions, saturation-state schedule, convergence criteria,
representative-slice medoid calculation, and endpoint-mapping equations.}

\SISubsection{Reservoir-ready consistency checks}

Each reservoir-ready property file is checked for complete $6\times87$
coverage, finite positive permeability, valid porosity, monotonic capillary-
pressure curves, physical relative-permeability bounds, and matching
saturation endpoints. The companion stratigraphy file must have the same
geology identifier and compatible content hash. Additional checks verify the
permeability units, local-to-reservoir tensor transformation, saturation-
region lookup tables, realization provenance, and one-to-one mapping between
the sampling manifest and replayed properties. Export proceeds only after all
checks pass.

\todo{Tabulate every acceptance criterion and report validation coverage for
the complete production dataset.}
